% Page budget: 1.55 pages (header value; README ownership now adds PS2, optimiser and checkpoint selection) | Owns: augmented model, encoder, closed-loop identification problem, optimiser and selection, added-state initialisation (PS2); models M_U, M_PS2.
% WRITING GUIDE
% Purpose: Define the final augmentation architecture and explain how it is trained inside the known feedback loop.
% Start here: Current implementation decisions D-129, D-140, D-141, D-180, D-220, D-233, and D-237 in docs/decisions.md.
% Supporting material: Quinten's 18 September outline feedback, the Meeting-audit synthesis, Thesis-writeup/Documentation/TRAINING-DESIGN.md, and docs/writeup figures.
% Mandatory papers: Read Hoekstra's 2025 LFR augmentation paper, the 2026 first-principles augmentation paper, the encoder-initialization paper, and Kessels' Chapter 5 before writing this section.
% Research-plan anchor: Preserve Aspect 2's dynamic parallel augmentation objective; document the departure from open-loop training and earlier state routing.
% Current decisions: Separate physical and additional states, keep the controller outside the model, score the complete training window without burn-in, and use one network for the physical correction and added-state update.
% Choices to justify: Final reported routing, added-state dimension, ANN width and depth, zero initialization, encoder history, closed-loop objective, rollout length, full-window scoring, and validation horizon. Use the configuration of the reported thesis run, not a superseded routing experiment.
% Horizon check: Treat encoder history, scored training window, joint-estimation horizon, and free-run validation horizon separately; relate candidates to relevant stable modes and test sensitivity to slower dynamics and free-integrator accumulation.
% Implementation audit: Read scripts/gantry/gantry_interconnect_dynamic.py, gantry_dynamic/{model,config,controller,data,training}.py, and model_augmentation/fit_systems/{interconnect,closed_loop,pre_encoder}.py.
% Reference check: Verify sources for parallel LFR augmentation, encoder initialization, closed-loop state extension, and any claimed architectural precedent against the original paper or thesis.
% Cautions: Earlier routing plans are superseded; distinguish the truth-only hidden absorber from learned states; closed-loop fit alone does not prove autonomous plant recovery.
% Figures: Absorber schematic, author's updated architecture, and thesis-owned common-reference loops. No residual panel or horizon figure.
% Section is finished when: Every signal has a defined frame and timing, the RK4 boundary is explicit, and the described architecture matches the final code.
%
% SOURCE MAP
% Hoekstra et al. 2025 EJC: dynamic-parallel topology, additional states, encoder-based truncated simulation.
% Hoekstra et al. 2026 Automatica submission: extended LFR formulation, baseline-preserving model/encoder initialization, joint identification context.
% Hoekstra et al. 2026 encoder paper: reconstructability-based W^b initialization; it does not identify W^a or test dynamic augmentation.
% Kessels 2025 thesis: direct closed-loop rollout, controller equations, and reconstruction of the machine controller state for each window.
% Drenth 2025 thesis: LPV-specific augmentation precedent; not the unrestricted MLP structure used here.
% This project: gantry specialization, residual-loop implementation, exact routing, RK4 boundary, and verification.
%
% MUST ESTABLISH (agreed with Dirk, 2026-10-06; step 2a of the thesis-section skill)
% Contribution delivered: the gantry realisation of S-DP augmentation (CT self-scheduled LPV-LFR baseline
%   inside a DT parallel model, joint estimation in the ten combinations) and its identification under
%   feedback. Not new: the S-DP topology (Hoekstra), augmentation of self-scheduled LPV-LFR with added
%   states (Drenth Eq. 5.2), closed-loop training with a known controller (Kessels, inspiration only).
% III-A Augmented model
%  1. Dynamic, not static: static augmentation adds no states (EJC Sec. 2, "flexible modes"); refitting
%     theta_base adds none either. Argue from the general case, no absorber.
%  1b. The augmented model is Hoekstra's LFR augmentation structure (EJC Eq. 4, 2026 Eq. 6) with a FIXED
%     interconnection (EJC Sec. 3.2) realising S-DP; its baseline block is Section II's self-scheduled
%     LFR with Delta(Y) = Y I6: an LFR inside an LFR, which is what the section title means.
%     Well-posedness: acyclic graph (EJC Thm. 1, 2026 Cond. 6) plus Section II's M(Y) condition.
%  2. Parallel, not series. LEAD REASON (Dirk, 2026-10-06): the orthogonal-by-construction method that
%     Section IV builds on is defined for the additive (parallel) structure (Gyorok 2026 abstract, Sec. 1;
%     D-003), so parallel lets us use that existing work. Supporting: Hoekstra gives no universal rule,
%     the structure follows prior knowledge (2026 Sec. 1, 6.3); the baseline stays the explicit nominal
%     term (research plan Aspect 2); parallel needs only a feedforward network for baseline-preserving
%     init, series needs ResNets (EJC Sec. 4.3, 2026 Sec. 5.4, 6.3).
%  3. State augmentation only, h_aug = 0: output P^T q is exact kinematics (cf. Kessels Remark 5.3). TODO reason.
%  4. Correction on all physical rows and added states: S-DP definition (EJC), X and Y reachable (D-103);
%     marginal-mode risk on free axes stated. TODO: why correct position rows rather than keep kinematics.
%     Context: Hoekstra 2026 Sec. 7 (F1Tenth) detaches the free integrators and learns only the
%     input-to-velocity map because the full state is measured; here only positions are measured.
%     One network for f_aug and g_aug is an implementation statement (D-180), no scientific claim.
%  5. DT correction after the CT RK4 step (input held, M(Y) per stage, affine normalisation). TODO reason
%     versus force injection in the CT derivative.
%  6. Unrestricted nonlinear augmentation (depends on Y through the state) versus LPV-structured
%     augmentation (Drenth): departure from the research plan's "in LPV-LFR form". TODO reason. Acyclic
%     routing adds no instantaneous loops; baseline well-posedness is Section II's.
%  7. Zero output layer reproduces the baseline from the same physical initial state (Hoekstra 2026
%     Eq. 29, D-237); the split is not unique (Sec. 5.2), handed to Section IV. Added states are memory
%     and model order, defined up to a state transformation, no physical coordinates.
% III-B Encoder
%  8. Encoder needed: only positions measured; truncated windows for cost and stability (Beintema).
%  9. W^b from the reconstructability map of a ZOH linearisation at an equilibrium (encoder paper
%     Eqs. 15-17, 30-32), not fitted (Hoekstra 2026 Eq. 30); W^a random. Claim only better starting
%     values and faster convergence, final accuracy comparable (encoder paper Sec. 4.4). The paper
%     assumes a stable baseline; ours is marginally stable (free X, Y), the map exists because the
%     positions make the linearisation observable. The paper treats co-estimating theta_base as out of
%     scope (footnote 1); we co-estimate. Precedent for linearising at zero: encoder paper Sec. 4.2
%     ("around the 0 point"), no reason given. TODOs: why Y = 0; disclose nominal parameters in W^b
%     versus detuned start; why model-based rather than data-based init (candidate: cost, Sec. 4.4);
%     Kaiming versus Xavier (todo only).
% 10. Encoder estimates the full state incl. measured positions (SUBNET). TODO: mention Kessels Remark 5.3?
% 11. Keep encoder history, scored window and validation horizon apart; values in Setup. The
%     joint-estimation sensitivity horizon belongs to Section IV.
% III-C Closed-loop identification
% 12. Known-controller, reference-driven output-error identification, inspired by Kessels (Eqs. 5.12,
%     5.13; Remark 5.6 open-loop attempt failed). Related to the indirect approach (Forssell and Ljung
%     Sec. 5.2) as a linear-theory motivation, not a guarantee; direct identification is not called
%     inherently biased. Free X and Y axes: the SUBNET consistency argument that Hoekstra inherits
%     (2026 Sec. 5.5) assumes an incrementally exponentially output stable data-generating system
%     (Beintema 2023 Condition 1); the open-loop gantry is not (free integrators), the loop with the
%     stabilising controller can be. Precondition, not a proof, for our nonlinear model. Same principle
%     as Kessels (closed-loop simulation of an unstable open-loop system), scoped to marginal stability.
% 13. Our formulation (ours, "we"): u_hat = u + C x^c + D(y - y_hat). Derivation by subtracting the
%     machine's controller equations, equal to the shared-r, shared-u_ff loop of the figure under a
%     compatible controller, signals and initial state (assumption; whether the data meet it is a Setup
%     or Results check: 4 kHz controller on recorded r - y versus recorded feedback force).
%     Main consequence: x^c_tau = 0 per window, no controller initialisation; consistent with the
%     encoder (model state and model controller both start on the machine's condition at tau). Contrast:
%     Kessels' direct form needs the machine controller state (Remark 5.4). Error before tau is not
%     carried in; the free-run validation does carry it.
% 14. Step order from no plant feedthrough: no algebraic loop, no added delay.
% 15. Every window sample scored (as Eq. 22); theta_base, theta_aug, eta joint, controller fixed.
%     Joint estimation framing (Dirk, 2026-10-06): the goal is to estimate theta_base jointly, in the ten
%     combinations. Obstacle from Jan's work: the split is not unique (2026 Sec. 5.2) and approximately
%     initialised parameters stay near their start (EJC Sec. 4.3, 2026 Sec. 6.3). Handover: Section IV
%     extends Jan's method with orthogonality so that the estimate of theta_base is UNIQUE (this
%     strength, not "recovers the true parameters"). State what interpretable means: the estimated
%     parameters keep their physical meaning; cite Gyorok 2026 Sec. 3 ("With non-unique theta_b, the
%     baseline model could lose its physical meaning"). Operational form only; the general definition
%     is the Introduction's (todo there). "Negation" and Kessels' definition (Sec. 5.3.1.3) belong to
%     Section IV; his Sec. 6.2 recommendation to the Introduction gap. Do NOT mention Jan's parameter regularisation anywhere.
%     Do NOT say explicitly that Section III's formulation is the unconstrained comparison arm; let
%     Section IV extend it naturally.
% Setup pattern for hyperparameters (2026 Sec. 7): honest basis per value (physical insight and
%     trial-and-error, line-search, inherited from earlier results).
% 16. Closed-loop fit can mask plant error; plant-domain checks and controller transfer in Results.
% Notation: u = P u_act as in Sections II and IV; no "(eom)" shorthand; controller matrices distinct from A_c(Y).
%
% PROPOSED MUST ESTABLISH (sandbox writer, cycle 04, 2026-10-09; NOT approved by Dirk; inferred from the README
%   ownership row III, source-map rows III-C and III-D, DOC/PS2-METHOD.md, DOC/THESIS-RESULTS.md and the code)
% III-C additions: (a) the identification problem is one display, cost subject to model, encoder and residual loop
%   (Hoekstra 2026 Eq. 22 pattern); (b) theta_base is trained in the coordinates xi of eq:mdelta_map, which keep
%   eq:lfr_admissibility at every iterate (D-229); (c) Adam then L-BFGS (Drenth Sec. 5), checkpoint selection by the
%   closed-loop free-run RMS on the validation records, polish kept only if it improves that score and no meter
%   regresses (D-171, D-172, closed_loop.py::validation_error, lbfgs_polish.py::decide); (d) model name M_U.
% III-D PS2 (D-234, D-236, D-240, D-243; FS/ps2_opening.py): 1. problem: at the zero start V gives W^a and the
%   added-state rows of W_L zero gradient; 2. S1 objective displayed (decoder on psi^a, target = current model's
%   closed-loop residual over n_p = n+1 samples, unit batch RMS), interleaved one step per Adam update of V;
%   3. S2 regression displayed (g_aug on frozen encoder pairs); 4. switching rules as the code states them (S1
%   plateau or cap; screen kappa > kappa_max skips S2; S2 plateau or cap); 5. S3 = V unchanged, decoder discarded;
%   6. adapted from Hefny 2015 Sec. 3 (two-stage regression) and Downey 2017 Sec. 4.2 (2SR then BPTT), with the
%   differences; 7. model name M_PS2. Left out: auxiliary learning rate, batch sizes, snapshot sizes, calibration
%   split rule, gradient-buffer handling, fall-back details, evidence (Results).
% Symbols changed against the reset draft to remove clashes: encoder parameters theta_enc (eta_0 is the m_Delta
%   offset), machine controller state chi (s is the m_Delta scale), stacked windows bold y, u (xi is the training
%   coordinate), PS2 horizon n_p (M is the inertia), decoder d_omega, unexplained fraction kappa.
\section{Dynamic LPV-LFR Augmentation}\label{sec:augmentation}
\ifdraft
\begin{itemize}
\item Given the baseline and closed-loop records, the goal is a model that adds the missing dynamics while $\theta_{\mathrm{base}}$ is estimated jointly; the section names its four subsections (style profile)
\item Contribution share: the second contribution of Section~I (dynamic parallel augmentation in a closed-loop, self-scheduled setting with added states), plus PS2 (MUST ESTABLISH; PROPOSED)
\end{itemize}
\fi

Given the baseline of Section~\ref{sec:lfr} and records measured under
feedback, our goal is a model that adds the dynamics the baseline omits while
$\theta_{\mathrm{base}}$ is estimated jointly.
To this end, Section~\ref{sec:aug_structure} defines the augmented model,
Section~\ref{sec:aug_encoder} the initial state of each training window,
Section~\ref{sec:aug_closed_loop} the closed-loop identification problem, and
Section~\ref{sec:aug_ps2} the initialisation of the added states.
Together they extend dynamic parallel LFR augmentation to a closed-loop,
self-scheduled state-space setting with learned additional states, the second
contribution of Section~I.
\todo{Missing: contribution wording. This section also delivers PS2 (Section~\ref{sec:aug_ps2}), which the Introduction's second contribution does not name. Candidate: ``\dots with learned additional states, initialised from data by a two-stage predictive-state regression'' (D-236, THESIS-RESULTS step~3).}

\subsection{Augmented model}\label{sec:aug_structure}
\ifdraft
\begin{itemize}
\item If the omitted effects are dynamic, the six baseline states cannot represent them; refitting $\theta_{\mathrm{base}}$ or a static augmentation adds no states (EJC Sec.~2)
\item The model is Hoekstra's LFR augmentation structure with a fixed interconnection realising S-DP; its baseline block is the LFR of Section~II; the graph is acyclic (EJC Sec.~3.2, Thm.~1; 2026 Table~1)
\item Realisation choices in one paragraph: parallel because OBC is defined for the additive structure (Gy\"or\"ok 2026); output map not learned (?); all physical rows corrected, $X$ and $Y$ meet no stiffness (?); no LPV structure, unlike Drenth Eq.~(5.2) (?)
\item $f_{\mathrm{base}}$ is one RK4 step with the input held and $M(Y)$ per stage, so it stays self-scheduled; the model runs in normalised coordinates, centred and with state statistics from the measured positions, unlike Hoekstra Eq.~(28) (code \texttt{compute\_normalization})
\item One tanh MLP supplies $f_{\mathrm{aug}}$ and $g_{\mathrm{aug}}$; a feedforward network suffices for a parallel structure (EJC Sec.~4.3; D-180, D-240)
\item A zero output layer reproduces the baseline (2026 Eq.~(29)); the added states have no physical coordinates
\item The baseline/network split is not unique (2026 Sec.~5.2, 6.3); interpretable means the estimated $\theta_{\mathrm{base}}$ keeps its physical meaning (Gy\"or\"ok 2026 Sec.~3); Section~IV makes the estimate unique
\end{itemize}
\fi

If the effects that the baseline omits are dynamic, its six states cannot
represent them.
Refitting $\theta_{\mathrm{base}}$ adds no states.
Neither does a static augmentation, which corrects the baseline without
extending its state~\cite[Sec.~2]{hoekstra2025lfr}.
The augmentation therefore needs states of its own.

We use the LFR-based augmentation structure of Hoekstra et al., in which a
fixed interconnection connects a baseline block and learning
blocks~\cite[Sec.~3.2]{hoekstra2025lfr}.
In our realisation, the baseline block is the self-scheduled LFR of
Section~\ref{sec:lfr}, so the augmented model is an LFR whose baseline block
is itself an LFR in $Y$.
The interconnection realises the dynamic-parallel
structure~\cite[Table~1]{hoekstra2026lfr}
\begin{equation}
\begin{aligned}
 \tilde x_{k+1}
 &=f_{\mathrm{base}}(\theta_{\mathrm{base}},\tilde x_k,u_k)
   +f_{\mathrm{aug}}(\theta_{\mathrm{aug}},\hat x_k,u_k),\\
 \bar x_{k+1}
 &=g_{\mathrm{aug}}(\theta_{\mathrm{aug}},\hat x_k,u_k),\\
 \hat y_k&=h_{\mathrm{base}}(\tilde x_k),
\end{aligned}
\label{eq:aug_transition}
\end{equation}
where $\tilde x_k=\col(q_k,\dot q_k)\in\R^6$ is the physical state of
\eqref{eq:lfr_G}, $\bar x_k\in\R^{n_{\bar x}}$ the added state,
$\hat x_k=\col(\tilde x_k,\bar x_k)$, $u_k=Pu_{\mathrm{act},k}\in\R^3$ the
generalised force of \eqref{eq:Ptransform},
$f_{\mathrm{base}}:\R^{10}\times\R^6\times\R^3\to\R^6$,
$f_{\mathrm{aug}}:\R^{n_\theta}\times\R^{6+n_{\bar x}}\times\R^3\to\R^6$ and
$g_{\mathrm{aug}}:\R^{n_\theta}\times\R^{6+n_{\bar x}}\times\R^3\to\R^{n_{\bar x}}$
the transition maps, $h_{\mathrm{base}}:\R^6\to\R^3$,
$h_{\mathrm{base}}(\tilde x)=P^\top q$, the output map to the actuator
positions $\qact$, and $\theta_{\mathrm{base}}\in\R^{10}$ from \eqref{eq:phi}
and $\theta_{\mathrm{aug}}\in\R^{n_\theta}$ the parameters.
The interconnection graph is acyclic, so it adds no algebraic loop to the
well-posed scheduling loop of Section~\ref{sec:lfr}~\cite[Thm.~1]{hoekstra2025lfr}.

The structure is parallel, augments the state only, and corrects every
physical row.
We choose the parallel interconnection because the orthogonal-by-construction
method that Section~\ref{sec:obc} extends is defined for this additive
structure~\cite[Abstract]{gyorok2026obc}.
The output map is not augmented.
As in the dynamic-parallel structure~\cite[Table~1]{hoekstra2026lfr},
$f_{\mathrm{aug}}$ writes all six physical rows, so the learned correction
reaches the $X$ and $Y$ axes.
On these axes it meets no stiffness in
\eqref{eq:damping_stiffness_matrices}, so a persistent correction accumulates
in position.
The network receives $Y$ through $\hat x$.
Unlike the augmented LPV-LFR of Drenth~\cite[Eq.~(5.2)]{drenth2025thesis}, it
has no LPV structure.
\todo{Missing: reason for state augmentation only ($h_{\mathrm{aug}}=0$). Candidate: $h_{\mathrm{base}}$ is the exact kinematic map of the measured positions, the case in which Kessels sees no need to augment the output equation~\cite[Remark~5.3]{kessels2025ai}; an output augmentation that depends on the input would make $\hat y$ depend on the feedback force that depends on $\hat y$~\cite[Remark~5.1]{kessels2025ai}.}
\todo{Missing: why the position rows are corrected rather than kept as the integrals of the velocities. Context: Hoekstra et al.\ detach the integrators of their vehicle and identify only the input-to-velocity map, with the full state measured~\cite[Secs.~7.2, 7.4]{hoekstra2026lfr}; here only positions are measured. D-103 requires only that $X$ and $Y$ are reachable.}
\todo{Missing: reason for an unrestricted network rather than an LPV-structured augmentation. The research plan (Aspect~2) proposed the interconnection in LPV-LFR form and D-017 is not superseded; the title may then also need a change.}
\todo{Missing: the header's research-plan anchor asks to document the departure from earlier state routing. Candidate: leave it out, since earlier routing is development history (thesis-section skill, reader test), and state only the final routing (D-222).}

$f_{\mathrm{base}}$ is one RK4 step of \eqref{eq:eom} over the sample period
$T_s$, with the input held over the step.
Each stage evaluates $M(Y)$ at its own state, so $f_{\mathrm{base}}$ remains
self-scheduled.
The model is trained in coordinates normalised with statistics of the
training records.
Hoekstra et al.\ scale without centring and take the state statistics from a
baseline simulation~\cite[Sec.~5.3]{hoekstra2026lfr}.
We centre every signal and take the state statistics from $q=P^{-\top}y$ and
its first difference.
With these statistics, as in~\cite[Eq.~(28)]{hoekstra2026lfr},
\begin{equation}
\begin{aligned}
 &\tilde x^{\mathrm n}=T_x(\tilde x-\mu_x),\quad
 u^{\mathrm n}=T_u(\uact-\mu_u),\quad
 y^{\mathrm n}=T_y(\qact-\mu_y),\\
 &f^{\mathrm n}_{\mathrm{base}}(\theta_{\mathrm{base}},\tilde x^{\mathrm n},u^{\mathrm n})
 =T_x\bigl(f_{\mathrm{base}}(\theta_{\mathrm{base}},T_x^{-1}\tilde x^{\mathrm n}+\mu_x,
 P(T_u^{-1}u^{\mathrm n}+\mu_u))-\mu_x\bigr),
\end{aligned}
\label{eq:aug_norm}
\end{equation}
where $y\in\R^3$ is the recorded output, $\mu_x\in\R^6$ and
$\mu_u,\mu_y\in\R^3$ are channel means, and $T_x$, $T_u$, $T_y$ are diagonal
with the inverse channel standard deviations.
\todo{Missing: reason for taking the state statistics from the measured positions rather than from a baseline simulation. Candidate: an input-driven baseline simulation drifts on the free $X$ and $Y$ axes (Section~\ref{sec:aug_closed_loop}) (own reasoning; D-119 removed the simulated-state source).}
\todo{Notation conflict with Section~\ref{sec:obc}: \eqref{eq:obc_tuples} defines the reference state with $T_x$ and $\mu_x$, so it is a normalised state, but \eqref{eq:obc_model} feeds it to $f_{\mathrm{base}}$ under the symbol $\tilde x$, which Section~\ref{sec:lfr} and this section use for the physical state. Candidate: write $f^{\mathrm n}_{\mathrm{base}}$ and $\tilde x^{\mathrm n}$ of \eqref{eq:aug_norm} in Section~\ref{sec:obc}.}

One multilayer perceptron with two tanh hidden layers supplies both learned
functions in these coordinates,
\begin{equation}
 \begin{bmatrix}
  T_xf_{\mathrm{aug}}(\theta_{\mathrm{aug}},\hat x,u)\\
  g_{\mathrm{aug}}(\theta_{\mathrm{aug}},\hat x,u)
 \end{bmatrix}
 =\phi_{\mathrm{aug}}(\theta_{\mathrm{aug}},\tilde x^{\mathrm n},\bar x,u^{\mathrm n}),
 \label{eq:aug_network}
\end{equation}
where $\phi_{\mathrm{aug}}:\R^{n_\theta}\times\R^{6+n_{\bar x}}\times\R^3\to\R^{6+n_{\bar x}}$
and $\theta_{\mathrm{aug}}$ are its weights and biases
(Fig.~\ref{fig:aug_structure}).
A feedforward network suffices because the baseline is augmented additively,
whereas a series structure needs a residual network to reproduce the baseline
at initialisation~\cite[Sec.~4.3]{hoekstra2025lfr}.

\begin{figure}[t]
 \centering
 \resizebox{\columnwidth}{!}{\input{figures/tex/augmentation_structure}}
 \caption{Augmented model. The physics step $f_{\mathrm{base}}$ and the
 network $\phi_{\mathrm{aug}}$ act in parallel on $\hat x_k$.
 $\phi_{\mathrm{aug}}$ supplies $f_{\mathrm{aug}}$ and $g_{\mathrm{aug}}$.
 The output map is not learned ($h_{\mathrm{aug}}=0$). The encoder
 $\psi$ sets $\hat x$ at each window start. During training, $u_k$ is the
 closed-loop input $\hat u_k$.
 \todo{No setpoint generator currently for the inputs. to cramped together. We could make this a two column figure.}
 }
 \label{fig:aug_structure}
\end{figure}

Training starts from the baseline.
The output layer of $\phi_{\mathrm{aug}}$ is initialised at zero,
\begin{equation}
 W_L=0,\quad b_L=0
 \quad\Longrightarrow\quad
 f_{\mathrm{aug}}=0,\quad g_{\mathrm{aug}}=0,
 \label{eq:aug_zero_init}
\end{equation}
where $W_L$ and $b_L$ are its weight and bias.
From the same physical initial state, the augmented model then reproduces the
baseline, as~\cite[Eq.~(29)]{hoekstra2026lfr} requires of an initialisation.
\todo{Missing: attribution of the all-zero output layer. It follows the authors' public dynamic-parallel code (D-237); \cite[Sec.~5.4.3]{hoekstra2026lfr} also allows a random linear part. Ask Jan which produced the reported S-DP results.}

Neither the added states nor the split between baseline and network are
fixed by the data.
An invertible change of coordinates of $\bar x$, composed into
$f_{\mathrm{aug}}$ and $g_{\mathrm{aug}}$, leaves the input-output behaviour
unchanged, so the added states have no physical coordinates.
$f_{\mathrm{aug}}$ can also learn changes that a different
$\theta_{\mathrm{base}}$ would produce, so the split is not
unique~\cite[Sec.~5.2]{hoekstra2026lfr}.
In the experiments of Hoekstra et al., the network learns dynamics that the
baseline could represent~\cite[Sec.~6.3]{hoekstra2026lfr}.
With a non-unique $\theta_{\mathrm{base}}$, the baseline can lose its physical
meaning~\cite[Sec.~3]{gyorok2026obc}.
We call the estimated combinations interpretable if they keep that meaning,
which requires a unique estimate.
Section~\ref{sec:obc} extends the method so that the estimate is unique.
\todo{Missing: the header Caution ``distinguish the truth-only hidden absorber from learned states''. The absorber is introduced only in Section~\ref{sec:setup}, so this section states the general fact (no physical coordinates). Candidate: one clause in Section~\ref{sec:setup} that the absorber states are not model states.}

\subsection{Encoder and initial state}\label{sec:aug_encoder}
\ifdraft
\begin{itemize}
\item Each window starts from a state of which only the positions are measured; an encoder estimates it, and short windows keep the cost independent of the record length and stabilise the optimisation (Beintema Secs.~2, 3)
\item The encoder is a linear map plus a nonlinear correction on a window ending at $\tau$ (encoder paper Eq.~(8), Eq.~(15)); unlike 2026 Eq.~(22e), the current sample is included
\item $W^a$ starts from a Xavier draw (2026 Eq.~(31); D-239)
\item $W^b$ is the reconstructability map of the baseline linearised at rest, $Y=0$, nominal parameters (encoder paper Eqs.~(16), (17)); computed, not fitted as in 2026 Eq.~(30)
\item The paper assumes a stable baseline; ours has two eigenvalues at one, but the positions make the linearisation observable (checked: rank of $O_n$ is 6)
\item The initialisation improves starting values and convergence speed, not final accuracy (encoder paper Sec.~4.4)
\item Open: why $Y=0$, nominal parameters versus a detuned start, model-based versus data-based initialisation (?)
\end{itemize}
\fi

Each training window starts from a state of which only the positions are
measured.
As in SUBNET~\cite{beintema2023subnet}, an encoder estimates this state from
the recorded input and output.
Simulating short windows from estimated states keeps the cost of a gradient
step independent of the record length and stabilises the
optimisation~\cite[Secs.~2, 3]{beintema2023subnet}.
The encoder is a linear map with a nonlinear
correction~\cite[Eq.~(8)]{hoekstra2026encoder},
\begin{equation}
\begin{aligned}
 \begin{bmatrix}T_x(\tilde x_{\tau|\tau}-\mu_x)\\ \bar x_{\tau|\tau}\end{bmatrix}
 &=\psi(\theta_{\mathrm{enc}},\mathbf y_\tau,\mathbf u_\tau)
 =\begin{bmatrix}W^b\\ W^a\end{bmatrix}
 \begin{bmatrix}\mathbf y_\tau\\ \mathbf u_\tau\end{bmatrix}
 +\tilde\psi(\mathbf y_\tau,\mathbf u_\tau),\\
 \mathbf y_\tau&=\col(y^{\mathrm n}_{\tau-n},\dots,y^{\mathrm n}_\tau),\quad
 \mathbf u_\tau=\col(u^{\mathrm n}_{\tau-n},\dots,u^{\mathrm n}_\tau),
\end{aligned}
\label{eq:aug_encoder}
\end{equation}
where $\tau$ is the window start, $n$ the encoder history,
$\psi:\R^{n_{\mathrm{enc}}}\times\R^{3(n+1)}\times\R^{3(n+1)}\to\R^{6+n_{\bar x}}$,
$\tilde\psi$ a tanh multilayer perceptron with the same domain and codomain
whose output layer starts at zero,
and $\theta_{\mathrm{enc}}\in\R^{n_{\mathrm{enc}}}$ collects $W^b$, $W^a$ and
the weights of $\tilde\psi$.
The window ends at $\tau$, as the reconstructability map below requires,
whereas the encoder of~\cite[Eq.~(22e)]{hoekstra2026lfr} stops at $\tau-1$.
\todo{Missing: whether to mention the alternative of initialising the positions from the measurement~\cite[Remark~5.3]{kessels2025ai}.}
\todo{Align the figure's past-only encoder-history labels with
\eqref{eq:aug_encoder}, which also uses the current sample.}

The rows $W^a$ have no baseline counterpart and start from a Xavier draw, as
in~\cite[Eq.~(31)]{hoekstra2026lfr}.
The rows $W^b$ start from the baseline, following the model-based
initialisation of~\cite[Sec.~3.1]{hoekstra2026encoder}.
For zero input, every rest state with $\Theta=0$ is an equilibrium of
\eqref{eq:eom}, because $K$ vanishes on $X$ and $Y$.
We linearise the baseline at rest at $Y=0$ with the nominal parameters of
Appendix~\ref{app:params} and discretise it by zero-order hold over $T_s$.
In the coordinates of \eqref{eq:aug_norm}, this gives $(A_d,B_d,C_d,D_d)$,
whose reconstructability map~\cite[Eqs.~(16), (17)]{hoekstra2026encoder} is
\begin{equation}
 W^b=\begin{bmatrix}
  A_d^nO_n^{+} & \;r_n-A_d^nO_n^{+}T_n
 \end{bmatrix},
 \label{eq:aug_wb}
\end{equation}
where $O_n\in\R^{3(n+1)\times6}$ is the observability matrix of $(A_d,C_d)$
over $n+1$ samples, $(\cdot)^+$ the pseudoinverse,
$T_n\in\R^{3(n+1)\times3(n+1)}$ the Toeplitz matrix of the Markov parameters,
$r_n\in\R^{6\times3(n+1)}$ the input-to-state map over the same samples, and
the two blocks act on $\mathbf y_\tau$ and $\mathbf u_\tau$.
Unlike~\cite[Eq.~(30)]{hoekstra2026lfr}, which fits the baseline encoder to
simulated baseline states, we compute the map directly.

The map exists although the baseline is not stable.
The encoder paper assumes a stable baseline~\cite[Sec.~2]{hoekstra2026encoder},
whereas $A_d$ has two eigenvalues at one, from the free $X$ and $Y$ axes.
The measured positions still make $(A_d,C_d)$ observable, so $O_n$ has full
column rank.
The model-based initialisation gives better starting values and faster
convergence than a random encoder, at comparable final
accuracy~\cite[Sec.~4.4]{hoekstra2026encoder}.
\todo{Missing: why linearise at $Y=0$. The encoder paper also linearises ``around the 0 point''~\cite[Sec.~4.2]{hoekstra2026encoder} but gives no reason. Candidate: $Y=0$ is the centre of the arm, from which $Y$ is measured~\cite[Sec.~2.2]{garcia2013model}, so it is the middle of the stroke.}
\todo{Missing: disclosure. $W^b$ uses the nominal parameters while $\theta_{\mathrm{base}}$ may start detuned (Section~\ref{sec:setup}) and is co-estimated, which the encoder paper leaves out of scope~\cite[footnote~1]{hoekstra2026encoder}. Decide whether to state that the encoder then starts inconsistent with $f_{\mathrm{base}}$.}
\todo{Missing: why the model-based rather than the data-based initialisation, which fits $W^b$ to forward-simulated baseline states~\cite[Sec.~3.2]{hoekstra2026encoder}. Candidate: that simulation drifts on the free $X$ and $Y$ axes (Section~\ref{sec:aug_closed_loop}), and the model-based map costs one pseudoinverse~\cite[Sec.~4.4]{hoekstra2026encoder} (own reasoning).}
\todo{VERIFY bib: \texttt{hoekstra2026encoder} = arXiv:2602.13108v1
(13 Feb 2026), Hoekstra, Gy\"or\"ok, T\'oth, Schoukens; matches PDF p.~1.
Check for a published (IFAC) version before submission.}

\subsection{Closed-loop identification}\label{sec:aug_closed_loop}
\ifdraft
\begin{itemize}
\item The records are measured under feedback and the model must predict the machine under feedback, so each window is simulated in closed loop with the known controller (Kessels Eqs.~(5.12), (5.13)), not driven by the recorded input (2026 Eq.~(22)); input-driven, the free axes accumulate position error
\item The consistency argument assumes an incrementally exponentially output stable system (2026 Sec.~5.5, Beintema Cond.~1), which the loop can meet and the open-loop gantry cannot; for linear systems this is the indirect approach (Forssell and Ljung); precondition, not proof
\item Machine and model loops share $r$, $u^{\mathrm{ff}}$ and the controller; subtracting them gives our residual form, exact for a compatible controller and signals (check in Setup)
\item $x^{c}_\tau=0$ needs no controller-state reconstruction, unlike Kessels Remark~5.4; no plant feedthrough fixes the step order
\item The identification problem: $V$ over every window sample subject to model, encoder and residual loop; $\theta_{\mathrm{base}}$ joint, unlike Kessels Eq.~(5.12)
\item $\theta_{\mathrm{base}}$ is trained in coordinates $\xi$: nine logarithmic, $m_\Delta$ through \eqref{eq:mdelta_map}, which keeps \eqref{eq:lfr_admissibility} at every iterate (this work; D-229)
\item Adam then L-BFGS (Drenth Sec.~5); selection by the closed-loop free-run RMS on the validation records; polish kept only if it improves and no meter regresses (D-171, D-172)
\item The identified model is $\mathcal M_{\mathrm U}$; a closed-loop fit can hide plant error (Results)
\end{itemize}
\fi

The records are measured under feedback, and the model must predict the
machine under feedback.
We therefore simulate each window in closed loop with the known controller of
its record, as Kessels does for augmented
models~\cite[Eqs.~(5.12), (5.13)]{kessels2025ai}.
Hoekstra et al.\ instead drive the model by the recorded
input~\cite[Eq.~(22)]{hoekstra2026lfr}.
In an input-driven simulation, a velocity error on the free $X$ and $Y$ axes
of \eqref{eq:damping_stiffness_matrices} accumulates in position.

Closing the loop also addresses the stability assumption behind consistency.
The consistency of the identification method~\cite[Sec.~5.5]{hoekstra2026lfr}
rests on that of SUBNET, which assumes an incrementally exponentially output
stable data-generating system~\cite[Cond.~1]{beintema2023subnet}.
A gantry with free axes does not meet this condition in open loop, whereas
the loop with its stabilising controller can.
For linear systems, this is the indirect approach, in which an output-error
model remains consistent because the reference is uncorrelated with the
noise~\cite[Sec.~5.2.1]{forssell1999revisited}.
For our nonlinear model, both arguments are preconditions, not a proof of
consistency.
\todo{Missing: evidence on the thesis data that open-loop rollouts accumulate position error on $X$ and $Y$ (Section~\ref{sec:results} bullet, THESIS-RESULTS Sec.~5).}

Both loops are driven by the same reference $r$ and feedforward
$u^{\mathrm{ff}}$ through the same controller, and differ only in the output
they feed back (Fig.~\ref{fig:aug_closed_loop}).
The controller of each record is linear, time invariant and known.
With machine controller state $\chi$ and model controller state $\hat\chi$,
the loops are
\begin{equation}
\begin{aligned}
 \chi_{k+1}&=A_{\mathrm{fb}}\chi_k+B_{\mathrm{fb}}(r_k-y_k),\\
 u_{\mathrm{act},k}&=C_{\mathrm{fb}}\chi_k+D_{\mathrm{fb}}(r_k-y_k)+u^{\mathrm{ff}}_k,\\
 \hat\chi_{k+1}&=A_{\mathrm{fb}}\hat\chi_k+B_{\mathrm{fb}}(r_k-\hat y_k),\\
 \hat u_{\mathrm{act},k}&=C_{\mathrm{fb}}\hat\chi_k+D_{\mathrm{fb}}(r_k-\hat y_k)+u^{\mathrm{ff}}_k,
\end{aligned}
\label{eq:aug_direct_loop}
\end{equation}
where $\chi_k,\hat\chi_k\in\R^{n_c}$, the matrices
$(A_{\mathrm{fb}},B_{\mathrm{fb}},C_{\mathrm{fb}},D_{\mathrm{fb}})$ map the
actuator position error to actuator forces, and $u_{\mathrm{act},k}$ is the
recorded input.

We implement the model loop by subtracting the machine's equations from the
model's.
This removes $r$ and $u^{\mathrm{ff}}$ and gives
\begin{equation}
\begin{aligned}
 x^{c}_{k+1}&=A_{\mathrm{fb}}x^{c}_k+B_{\mathrm{fb}}(y_k-\hat y_k),\qquad x^{c}_\tau=0,\\
 \hat u_{\mathrm{act},k}&=u_{\mathrm{act},k}+C_{\mathrm{fb}}x^{c}_k+D_{\mathrm{fb}}(y_k-\hat y_k),
\end{aligned}
\label{eq:aug_loop}
\end{equation}
where $x^{c}=\hat\chi-\chi$ is the difference of the controller states.
The two forms are equal when the recorded input was produced by this
controller from the same output samples.
\todo{Check for Setup: the training loop uses the controller re-discretised at the model rate, while the records were generated at a higher rate. Run the re-discretised controller on the recorded $r-y$ and compare its feedback force with the recorded one. A rollout with $\hat y=y$ gives $\hat u=u$ by construction and checks nothing.}

\begin{figure}[t]
 \centering
 \includegraphics[width=\columnwidth]{closed_loop_same_reference.pdf}
 \caption{Recorded plant loop (top) and augmented-model loop (bottom),
 driven by one shared reference $r_k$ and feedforward $u_k^{\mathrm{ff}}$,
 with the same feedback controller. The model output is evaluated
 before the input advances its state to the next sample.}
 \label{fig:aug_closed_loop}
\end{figure}

The residual form needs no controller initialisation.
Because $x^{c}$ is a difference of controller states, $x^{c}_\tau=0$ places
the model's controller in the machine's controller state at every window
start.
The encoder does the same for the model state.
Kessels' direct form instead needs the machine's controller state at each
window start, read from the machine or reconstructed from the reference, the
output and the controller~\cite[Remark~5.4]{kessels2025ai}.
The augmented model has no feedthrough, so each step evaluates $\hat y_k$
from $\tilde x_k$, then $\hat u_{\mathrm{act},k}$ from \eqref{eq:aug_loop},
and then advances both states.
The controller feedthrough $D_{\mathrm{fb}}$ therefore closes no algebraic
loop.

As in~\cite[Eq.~(22)]{hoekstra2026lfr}, the criterion scores every sample of
every window.
Over the window starts $\mathcal T$ and the window length $n_f$, the
identification problem is
\begin{equation}
 \min_{\vartheta}\;
 V(\vartheta)=\frac{1}{|\mathcal T|\,n_f\,n_y}
 \sum_{\tau\in\mathcal T}\sum_{\ell=0}^{n_f-1}
 \norm{y^{\mathrm n}_{\tau+\ell}-\hat y^{\mathrm n}_{\tau+\ell|\tau}}_2^2
 \label{eq:aug_loss}
\end{equation}
subject to \eqref{eq:aug_transition} with
$u_k=\hat u_k=P\hat u_{\mathrm{act},k}$, \eqref{eq:aug_encoder} and
\eqref{eq:aug_loop}, where $\vartheta=(\xi,\theta_{\mathrm{aug}},\theta_{\mathrm{enc}})$,
$\xi\in\R^{10}$ are the training coordinates of $\theta_{\mathrm{base}}$
below, $\hat y^{\mathrm n}$ is the model output normalised as in
\eqref{eq:aug_norm}, and $n_y=3$.
Unlike Kessels' criterion, which keeps the first-principles parameters
fixed~\cite[Eq.~(5.12)]{kessels2025ai}, $V$ is minimised over
$\theta_{\mathrm{base}}$ jointly with the network and the encoder.

$\theta_{\mathrm{base}}$ is not trained directly, because an unconstrained
update can violate \eqref{eq:lfr_admissibility}.
We train nine combinations in logarithmic coordinates and $m_\Delta$ in a
bounded one,
\begin{equation}
 \theta_{\mathrm{base},j}=\theta^0_{\mathrm{base},j}\,e^{\xi_j},\qquad
 m_\Delta=\rho\,\frac{2}{L_b}\sqrt{m_\Sigma J_{\mathrm{eff}}}\,\tanh(\eta_0+s\,\xi_\Delta),
 \label{eq:mdelta_map}
\end{equation}
where $j$ runs over the nine entries of \eqref{eq:phi} other than $m_\Delta$,
$\xi_\Delta$ is the entry of $\xi$ for $m_\Delta$, $\theta^0_{\mathrm{base}}$
the start, $\rho\in(0,1)$ a margin, and $\eta_0$ and $s$ make $\xi=0$ give
$\theta^0_{\mathrm{base}}$ (Appendix~\ref{app:lfr-wellposed}).
Since $\rho\lvert\tanh(\cdot)\rvert<1$, every $\xi$ satisfies
\eqref{eq:lfr_admissibility}, so $M(Y)\succ0$ and the baseline stays well
posed at every iterate without a clamp or a penalty.

We minimise $V$ with Adam on minibatches of windows and then refine the
result with L-BFGS, as Drenth et al.\ do~\cite[Sec.~5]{drenth2025lpvlfr}.
Checkpoints are selected by a closed-loop free run on the validation records.
Each record is simulated whole from the encoder state at its start, with
$x^{c}=0$ only there.
The selection score is the record-length-weighted mean of the per-record RMS
output error in metres.
The L-BFGS result is kept only if it lowers this score and no monitored
quantity regresses.
The values of $n$, $n_f$, $n_{\bar x}$, the network size and the optimiser
settings are given in Section~\ref{sec:setup}.
\todo{Missing: the L-BFGS acceptance monitors the combination error against the true parameters (\texttt{GD/training.py::\_gantry\_meters}, \texttt{combo\_err}), an oracle inside the training procedure (sources.md trap~6). Candidate: drop that meter from the acceptance rule for the reported runs, or state it as a simulation-only safeguard in Section~\ref{sec:setup}.}

We call the model \eqref{eq:aug_transition} with encoder
\eqref{eq:aug_encoder}, started from \eqref{eq:aug_zero_init} and identified
by \eqref{eq:aug_loss}, $\mathcal M_{\mathrm U}$.
A good closed-loop fit does not by itself show that the model is correct in
open loop, because the controller acts on the output error and can mask plant
error.
Section~\ref{sec:results} therefore also evaluates the model under a changed
controller.
\todo{Missing: model notation. Candidate: $\mathcal M_{\mathrm U}$ and $\mathcal M_{\mathrm{PS2}}$, since $M$ is the inertia matrix (README open conflict on model names, as for $\mathcal M_{\mathrm{LPV}}$ in Section~\ref{sec:lfr}).}

\subsection{Initialisation of the Added States}\label{sec:aug_ps2}
\ifdraft
\begin{itemize}
\item At the zero start, $V$ gives $W^a$ and the added-state rows of $W_L$ no gradient: the added states start without influence (PS2-METHOD Sec.~2.2)
\item PS2 adapts two-stage regression (Hefny Sec.~3) and refines its result by the simulation criterion (Downey Sec.~4.2) (this work, adapted)
\item S1: the encoder's added-state outputs and a temporary decoder predict the current model's closed-loop residual over $n_{\mathrm p}=n+1$ samples; one S1 step per Adam update of $V$ (\texttt{ps2\_opening.py})
\item S1 stops on a plateau of the unexplained fraction or a cap; above $\kappa_{\max}$, S2 is skipped (\texttt{PS2Spec})
\item S2: $g_{\mathrm{aug}}$ is fitted once to propagate the frozen encoder states by one sample; stops on a plateau or a cap
\item S3: the decoder is discarded and $V$ continues unchanged; the result is $\mathcal M_{\mathrm{PS2}}$
\end{itemize}
\fi

At the zero start \eqref{eq:aug_zero_init}, the output does not depend on
$\bar x$, so the gradient of $V$ with respect to $W^a$ and to the added-state
rows of $W_L$ is zero.
We therefore initialise the added states from data by PS2, two supervised
stages adapted from the two-stage regression of Hefny et
al.~\cite[Sec.~3]{hefny2015supervised}.

In the first stage, S1, the added-state outputs of the encoder learn to
predict the current model's closed-loop output residual over the next
$n_{\mathrm p}=n+1$ samples, the length of the encoder window, through a
temporary decoder $d_\omega$,
\begin{equation}
\begin{aligned}
 &\min_{\omega,\,\theta^{\mathrm a}_{\mathrm{enc}}}\;
 \sum_{\tau\in\mathcal B}
 \norm{d_\omega\bigl(\psi^{\mathrm a}(\theta_{\mathrm{enc}},\mathbf y_\tau,\mathbf u_\tau)\bigr)
 -\mathbf e_\tau}_2^2,\\
 &\mathbf e_\tau=\frac{\col\bigl(y^{\mathrm n}_{\tau+\ell}-\hat y^{\mathrm n}_{\tau+\ell|\tau}\bigr)_{\ell=0}^{n_{\mathrm p}-1}}
 {\operatorname{rms}_{\mathcal B}},
\end{aligned}
\label{eq:ps2_s1}
\end{equation}
where $\psi^{\mathrm a}$ are the added-state rows of \eqref{eq:aug_encoder},
$\theta^{\mathrm a}_{\mathrm{enc}}$ the entries of $\theta_{\mathrm{enc}}$
that feed only them ($W^a$ and the added-state rows of the output layer of
$\tilde\psi$), $d_\omega:\R^{n_{\bar x}}\to\R^{3n_{\mathrm p}}$ a one-layer
tanh network with parameters $\omega\in\R^{n_\omega}$ and zero output layer,
$\mathcal B$ a minibatch of window starts, $\hat y^{\mathrm n}_{\tau+\ell|\tau}$
the output of the current model, \eqref{eq:aug_transition} with
\eqref{eq:aug_loop}, from the encoder state without gradient, and
$\operatorname{rms}_{\mathcal B}$ the RMS of the unscaled residuals over the
batch.
Unlike the first stage of~\cite[Sec.~3]{hefny2015supervised}, which regresses
future features on the history directly, the prediction passes through the
$n_{\bar x}$-dimensional bottleneck $\psi^{\mathrm a}$.

S1 runs alongside the first updates of $V$: after every Adam update of $V$,
one S1 step updates $\omega$ and $\theta^{\mathrm a}_{\mathrm{enc}}$ with a
separate optimiser.
Its progress is the fraction $\kappa$ of $\sum_\tau\norm{\mathbf e_\tau}^2$
that $d_\omega$ leaves unexplained on windows of training records withheld
from the fit.
S1 ends when $\kappa$ stops decreasing, or after a maximum number of updates.
If $\kappa>\kappa_{\max}$ at the end of S1 or at an earlier screening update,
S2 is skipped.

In the second stage, S2, the deployed transition $g_{\mathrm{aug}}$ is fitted
once to propagate the encoder's added states by one sample,
\begin{equation}
 \min_{\theta_{\mathrm{aug}}}\;\sum_{i}
 \norm{g_{\mathrm{aug}}(\theta_{\mathrm{aug}},\hat x_i,u_{\tau_i})
 -\psi^{\mathrm a}(\theta_{\mathrm{enc}},\mathbf y_{\tau_i+1},\mathbf u_{\tau_i+1})}_2^2,
 \label{eq:ps2_s2}
\end{equation}
where $\hat x_i$ is the encoder state at a window start $\tau_i$ from the
records of the S1 fit, $u_{\tau_i}$ the recorded input, and all encoder outputs are computed once with
$\theta_{\mathrm{enc}}$ after S1 and then held fixed.
Unlike the second stage of~\cite[Sec.~3]{hefny2015supervised}, which regresses
a linear operator, the regressed map is the deployed nonlinear transition.

S2 stops when the same error on pairs from the withheld records stops
decreasing, or after a maximum number of steps.
The decoder is then discarded, and the minimisation of $V$ continues with the
optimiser and selection of Section~\ref{sec:aug_closed_loop}.
As in~\cite[Sec.~4.2]{downey2017psrnn}, the two-stage result thus initialises
a model that the simulation criterion refines.
We call the result $\mathcal M_{\mathrm{PS2}}$.
It differs from $\mathcal M_{\mathrm U}$ only by the S1 and S2 updates of
$\theta_{\mathrm{enc}}$ and $\theta_{\mathrm{aug}}$.
